{\footnotesize
\setlength{\tabcolsep}{4pt}
\begin{threeparttable}
\begin{tabular}{llrrrrrr}
\toprule
 & & & & \multicolumn{2}{c}{Monthly rate (\%)} & \multicolumn{2}{c}{Reemployed (\%)} \\
\cmidrule(lr){5-6}\cmidrule(lr){7-8}
Tenure & Group & Dismissed & \% of group & Months 1 to 5 & Months 6 to 8 & 6 m & 12 m \\
\midrule
Under 6 months & High (G1) & 2,728,702 & 16.6 & 8.17 & 6.68 & 39.4 & 56.4 \\
 & Medium (G2) & 4,187,646 & 21.5 & 7.18 & 6.41 & 35.8 & 53.0 \\
 & Low (G3) & 3,095,827 & 17.8 & 7.70 & 6.96 & 37.9 & 55.6 \\
\addlinespace[3pt]
24 months or more & High (G1) & 6,567,675 & 39.8 & 3.52 & 7.05 & 22.4 & 44.3 \\
 & Medium (G2) & 6,051,629 & 31.1 & 3.36 & 7.09 & 21.7 & 44.3 \\
 & Low (G3) & 6,489,698 & 37.2 & 2.92 & 7.92 & 20.6 & 45.3 \\
\bottomrule
\end{tabular}
\begin{tablenotes}[flushleft]\footnotesize
\item \textit{Note:} Six cohorts pooled. The monthly rate is the share of workers without a new job at the start of the month who are reemployed in that month, as a simple average over the months indicated. Workers with less than 6 months of tenure are not eligible for unemployment insurance; those with 24 months or more are entitled to five installments. Source: RAIS, authors' calculations.
\end{tablenotes}
\end{threeparttable}
}
